\documentclass[preprint,12pt]{elsarticle}
\usepackage[T1]{fontenc}
\usepackage[utf8]{inputenc}
\usepackage{lmodern}
\usepackage{microtype}
\usepackage{booktabs}
\usepackage{array}
\usepackage{tabularx}
\usepackage{longtable}
\usepackage{graphicx}
\usepackage{amsmath}
\usepackage{xurl}
\usepackage[hidelinks]{hyperref}
\newcolumntype{L}[1]{>{\raggedright\arraybackslash}p{#1}}
\setlength{\emergencystretch}{3em}
\makeatletter
\setlength{\@fptop}{0pt}
\makeatother
\setcounter{topnumber}{5}
\setcounter{bottomnumber}{5}
\setcounter{totalnumber}{10}
\renewcommand{\topfraction}{0.95}
\renewcommand{\bottomfraction}{0.90}
\renewcommand{\textfraction}{0.05}
\renewcommand{\floatpagefraction}{0.70}
\journal{Smart Agricultural Technology}
\setcounter{secnumdepth}{0}
% Author metadata supplied by the author team on 2026-08-11.
\newif\ifAuthorMetadataReady
\AuthorMetadataReadytrue

\newcommand{\CorrespondingAuthor}{Zhuo Chen}
\newcommand{\CorrespondingEmail}{zhuoc@chalmers.se}
\newcommand{\CorrespondingAffiliation}{Chalmers University of Technology, SE-412 96 Gothenburg, Sweden}

\newcommand{\AuthorFrontmatter}{%
  \def\thefnote{$\dagger$}%
  \author[aff-a]{Shuhao Liu\fnref{equal}}
  \ead{2420610137@stu.hrbust.edu.cn}
  \author[aff-b]{Zhuo Chen\fnref{equal}\corref{corresponding}}
  \ead{zhuoc@chalmers.se}
  \author[aff-c]{Zhi Ling}
  \ead{zhi.ling@kiwiar.com}
  \author[aff-c]{Yu Yan}
  \ead{yu.yan@kiwiar.com}
  \author[aff-a]{Jie Liu}
  \ead{liujie@hrbust.edu.cn}
  \author[aff-a]{Qiuxue Wu}
  \ead{2420610140@stu.hrbust.edu.cn}
  \author[aff-c]{Ziyi Kuang}
  \ead{ziyi.kuang@kiwiar.com}
  \address[aff-a]{Harbin University of Science and Technology, No. 52 Xuefu Road, Nangang District, Harbin 150080, Heilongjiang, China}
  \address[aff-b]{Chalmers University of Technology, SE-412 96 Gothenburg, Sweden}
  \address[aff-c]{Kiwiar Co., Ltd., Singapore}
  \fntext[equal]{These authors contributed equally to this work.}
  \cortext[corresponding]{Corresponding author.}
}

% These statements require separate confirmation from the author team.
\newif\ifFundingStatementReady
\FundingStatementReadyfalse
\newcommand{\FundingStatement}{}
\newif\ifCompetingInterestsReady
\CompetingInterestsReadyfalse
\newcommand{\CompetingInterests}{}
\newif\ifCRediTStatementReady
\CRediTStatementReadyfalse
\newcommand{\CRediTStatement}{}
\newcommand{\RepositoryDOI}{}


\begin{document}
\raggedbottom
\begin{frontmatter}
\title{Supplementary material: Support-Aware Review Queues for Crop--Weed Vision: A Multitemporal Single-Field Case Study}
\ifAuthorMetadataReady
\AuthorFrontmatter
\else
\author{Anonymous manuscript}
\fi
\end{frontmatter}

\section{S1. Public-data lineage and official spatial protocol}

This study uses two public datasets. SugarBeets2016 supplies source-specialist supervision, and WeedsGalore supplies target development and evaluation data. CropAndWeed and PhenoBench formed preliminary data-lineage audits and remain outside the reported experiments.

\begin{table}[htbp]
\centering
\caption{Official WeedsGalore split by acquisition date. Instance counts are from released raster enumeration.}
\small
\resizebox{\textwidth}{!}{%
\begin{tabular}{lrrrrr}
\toprule
Date & Train tiles & Validation tiles & Test tiles & Test crop & Test weed \\
\midrule
25 May 2023 & 32 & 8 & 8 & 122 & 392 \\
30 May 2023 & 32 & 8 & 8 & 55 & 440 \\
6 June 2023 & 32 & 8 & 8 & 68 & 726 \\
15 June 2023 & 8 & 2 & 2 & 11 & 154 \\
\midrule
Total & 104 & 26 & 26 & 256 & 1,712 \\
\bottomrule
\end{tabular}}
\end{table}

The public paper describes one field, 48 recorded locations, and an 8/2/2 spatial-patch allocation. The release supplies official split labels and capture indices. The manifest preserves this released spatial structure. Across all 156 rasters, enumeration yields 2,168 crop and 10,029 weed instances; the dataset paper reports 2,169 and 10,031. The consistency audit found zero role disagreements between instance and semantic rasters and zero excluded identifiers. Archive and manifest SHA-256 checksums are recorded in the reproducibility package.

\section{S2. Annotation and information-flow ledger}

\begin{table}[htbp]
\centering
\caption{Supervision ledger. ``Available'' distinguishes labels present in the public data from labels actually read by a stage.}
\small
\begin{tabularx}{\textwidth}{p{3.0cm}p{2.5cm}p{3.2cm}X}
\toprule
Information & Available & Read by & Purpose \\
\midrule
SugarBeets pixel masks & Public source & Specialist training & Train foreground and semantic specialists \\
Target train dense masks & 104 tiles & Proposal audit; label derivation & Select the commissioned proposal setting and derive candidate roles \\
Target validation dense masks & 26 tiles & Model/queue selection & Select $C$, fine-tuning epoch, semantic epoch, threshold or area, or queue-score threshold \\
Target test dense masks & 26 tiles & Offline scorer after declared outputs & Compute proposal, queue, IoU, and crop-overlap metrics \\
300 mask-derived training-candidate labels/draw & Simulated reveal from train pool & Label-budget ranker & Compare prospective acquisition policies \\
Proposal-generation inputs & RGB and source model outputs & Proposal generator & Create connected-component candidates \\
\bottomrule
\end{tabularx}
\end{table}

The full-supervision linear rankers use 9,887 official-training and 2,436 validation candidate labels before refitting. The equal-target-mask-access semantic baseline trains on all 104 target-training dense masks. The nominal 300-label experiment restricts only the sampled training candidates; all 2,436 labelled validation candidates remain available for model selection. Candidate labels are deterministically derived from dense masks: candidates satisfying the role-qualified weed contract are positive and all remaining candidates are negative. Prospective uniform and DINOv2 coreset strategies select training IDs before revealing labels. Retrospective class stratification uses labels during selection and serves as a retrospective reference.

The core configuration comprises the split, proposal grid, matching contract, feature definitions, and discrimination-selected rankers archived before final official-test scoring. Target-semantic, repeated-seed, quota, exact-mask, per-tile, crop-severity, and date/patch analyses were specified after the core test had been accessed and are exploratory. Training and validation data selected regularization, thresholds, epochs, and candidate geometry. Each test policy formed its queue from images, scores, tile identities, and declared group labels before the offline scorer accessed test masks.

\subsection*{Source specialist training and source-domain diagnostics}

Both source specialists used a ResNet-18 U-Net initialized from torchvision ImageNet-1K V1 weights, ImageNet RGB normalization, AdamW, bfloat16 autocasting, gradient clipping at 5, and the 7,050-image SugarBeets source training partition. The foreground specialist used 32,768 deterministic $128\times128$ patches, BCE with positive weight 30 plus soft Dice, and 10 epochs. The three-class semantic specialist used 16,384 role-balanced patches, weighted cross-entropy $(0.2,1,1)$ plus plant-role Dice, and 8 epochs. Training used deterministic patches, and the final epoch supplied each checkpoint. The 901-image source validation partition subsequently provided operating-point diagnostics.

On the later source validation diagnostic, the foreground model's maximum scanned pixel IoU was 0.3821 at threshold 0.99 (precision 0.5788, recall 0.5293); threshold 0.70 gave precision 0.0479 and recall 0.9631. The semantic model had raw three-class macro IoU 0.4425. Its independently selected role threshold produced balanced accuracy 0.7224, crop recall 0.5524, and weed recall 0.8924. These metrics document source behavior and checkpoint provenance; target queue conclusions use the official WeedsGalore split.

\section{S3. Proposal selection, definitions, and burden}

For truth instance $I$, candidate $C$, labeled pixels $L$, and pixels of role $R$, the audit records
\[
q_I=|I\cap C|/|I|,\qquad q_L=|C\cap L|/|C|,\qquad q_R=|C\cap R|/|C\cap L|.
\]
Spatial recall requires $q_I\ge0.50$. Role-qualified recall additionally requires $q_L\ge0.50$, $q_R\ge0.90$, and a role match. Queued recall restricts candidates to the selected queue. These fixed engineering thresholds define offline candidate accounting rather than agronomic treatment or actuation; the evaluated 27-cell sensitivity grid varies all three values.

\begin{table}[htbp]
\centering
\caption{Official-training proposal grid. Feasible settings have training-set mean burden at most 100 candidates/image, with tile-level counts unrestricted.}
\small
\resizebox{\textwidth}{!}{%
\begin{tabular}{lrrrr}
\toprule
Threshold / minimum area (px) & Candidates/image & Spatial weed recall & Qualified weed recall & Feasible \\
\midrule
0.30 / 8 & 281.35 & 0.7135 & 0.2406 & No \\
0.30 / 32 & 142.04 & 0.7000 & 0.2297 & No \\
0.50 / 8 & 259.49 & 0.6325 & 0.2822 & No \\
0.50 / 32 & 123.52 & 0.6190 & 0.2707 & No \\
0.70 / 8 & 207.48 & 0.5416 & 0.2979 & No \\
0.70 / 32 & 95.07 & 0.5271 & 0.2856 & Yes \\
0.80 / 8 & 168.63 & 0.4806 & 0.2959 & No \\
0.80 / 32 & 76.99 & 0.4672 & 0.2835 & Yes \\
0.90 / 8 & 115.50 & 0.3703 & 0.2421 & No \\
0.90 / 32 & 52.38 & 0.3587 & 0.2317 & Yes \\
0.95 / 8 & 75.71 & 0.2452 & 0.1685 & Yes \\
0.95 / 32 & 36.83 & 0.2366 & 0.1605 & Yes \\
\bottomrule
\end{tabular}}
\end{table}

The threshold 0.70, 32-pixel setting maximized training qualified recall among feasible settings. The deterministic tie order was qualified recall, lower mean burden, higher threshold, then larger minimum area. The earlier externally fixed setting used threshold 0.95 and the same minimum area and provides the shared-ranker-weight comparison.

\begin{table}[htbp]
\centering
\caption{Official-test proposal layers by date. Numerators are weed instances.}
\scriptsize
\resizebox{\textwidth}{!}{%
\begin{tabular}{llrrrr}
\toprule
Date & Setting & Weed $n$ & Spatial recalled & Qualified recalled & Candidates/image \\
\midrule
25 May & Original & 392 & 2 (0.0051) & 2 (0.0051) & 13.88 \\
25 May & Commissioned & 392 & 51 (0.1301) & 13 (0.0332) & 102.75 \\
30 May & Original & 440 & 106 (0.2409) & 84 (0.1909) & 35.75 \\
30 May & Commissioned & 440 & 263 (0.5977) & 121 (0.2750) & 119.38 \\
6 June & Original & 726 & 328 (0.4518) & 237 (0.3264) & 45.13 \\
6 June & Commissioned & 726 & 566 (0.7796) & 369 (0.5083) & 64.63 \\
15 June & Original & 154 & 73 (0.4740) & 7 (0.0455) & 22.00 \\
15 June & Commissioned & 154 & 98 (0.6364) & 15 (0.0974) & 28.50 \\
\midrule
Pooled & Original & 1,712 & 509 (0.2973) & 330 (0.1928) & 30.85 \\
Pooled & Commissioned & 1,712 & 978 (0.5713) & 518 (0.3026) & 90.42 \\
\bottomrule
\end{tabular}}
\end{table}

Candidate burden under the original setting had median 28.5, interquartile range 19.25--43.75, maximum 62, and 7/26 test images with fewer than 20 candidates. The commissioned setting had median 89, interquartile range 60--114.25, maximum 183, 10/26 images above 100, and no image with fewer than 20 candidates. Its 30 May mean was 130.38 on training and 119.38 on test. Candidate composition changed from 453 weed/148 crop/201 ambiguous (802 total) to 464/111/1,776 (2,351 total).

\begin{table}[htbp]
\centering
\caption{Split/merge sensitivity across the complete 156-tile release. ``Merge'' counts candidates overlapping at least two instances at the stated instance-coverage threshold.}
\small
\begin{tabular}{lrrr}
\toprule
Setting & Threshold & Split instances & Merge candidates \\
\midrule
Original & 0.10 & 224 & 735 \\
Original & 0.25 & 40 & 640 \\
Original & 0.50 & 0 & 473 \\
Commissioned & 0.10 & 154 & 1,136 \\
Commissioned & 0.25 & 26 & 1,065 \\
Commissioned & 0.50 & 0 & 938 \\
\bottomrule
\end{tabular}
\end{table}

\section{S4. Fixed-weight proposal comparison}

One all-candidate role-plus-local logistic ranker was fitted on commissioned train plus validation candidates. Its scaler and coefficients were then applied unchanged to both the original and commissioned test proposals. This isolates proposal-pool changes from ranker refitting.

\begin{table}[htbp]
\centering
\caption{Shared-ranker-weight comparison at $K=20$. Ranker weights are identical, but workloads differ because seven original-setting tiles supply fewer than 20 candidates. Crop frequency uses all 26 test tiles as denominator.}
\small
\resizebox{\textwidth}{!}{%
\begin{tabular}{lrrrrrr}
\toprule
Setting & Accepted & Precision & Spatial & Qualified & 1:1 qualified & Any crop \\
\midrule
Original & 471 & 276/471 & 321/1712 & 205/1712 & 166/1712 & 22/26 \\
Commissioned & 520 & 247/520 & 467/1712 & 273/1712 & 201/1712 & 23/26 \\
\bottomrule
\end{tabular}}
\end{table}

The complete $K=1,5,10,20,50$ curves and per-date rows are supplied in the reproducibility archive. All ties are resolved by lexicographically ascending candidate ID.

\section{S5. All-candidate ranking, per-date queue results, and ceilings}

\begin{table}[htbp]
\centering
\caption{Candidate-ranking support audit on the 2,351 official-test candidates.}
\scriptsize
\begin{tabular}{lrr}
\toprule
Model and fitting support & AUC & AP \\
\midrule
Source ResNet score, all candidates & 0.4947 & 0.2877 \\
Eligible role+local, all-candidate test & 0.4420 & 0.1640 \\
All-candidate role+local & 0.9049 & 0.6960 \\
All-candidate geometry & 0.8154 & 0.5418 \\
Eligible frozen DINOv2, all-candidate test & 0.6321 & 0.3816 \\
All-candidate frozen DINOv2 & 0.9630 & 0.8684 \\
All-candidate masked DINOv2 & 0.9665 & 0.8813 \\
\bottomrule
\end{tabular}
\end{table}

\begin{table}[htbp]
\centering
\caption{Component-masked DINOv2 queue at $K=20$ by date.}
\scriptsize
\resizebox{\textwidth}{!}{%
\begin{tabular}{lrrrrrr}
\toprule
Date & Accepted & Weed cand. & Precision & Spatial recall & Qualified recall & Any crop \\
\midrule
25 May & 160 & 13 & 0.0813 & 32/392 (0.0816) & 12/392 (0.0306) & 8/8 \\
30 May & 160 & 89 & 0.5563 & 128/440 (0.2909) & 84/440 (0.1909) & 2/8 \\
6 June & 160 & 150 & 0.9375 & 220/726 (0.3030) & 210/726 (0.2893) & 0/8 \\
15 June & 40 & 26 & 0.6500 & 79/154 (0.5130) & 13/154 (0.0844) & 1/2 \\
\midrule
Pooled & 520 & 278 & 0.5346 & 459/1712 (0.2681) & 319/1712 (0.1863) & 11/26 \\
\bottomrule
\end{tabular}}
\end{table}

For frozen DINOv2 at $K=20$, the role-qualified proposal ceiling is 518/1,712, the exact unconstrained oracle is 426/1,712, and the discrimination-selected masked queue is 319/1,712. The corresponding ranking gap is 0.0625. For a maximum crop-overlap scene frequency of 0.01, $\lfloor0.01\times26\rfloor=0$ scenes are allowed. Crop-excluding optimization yields 418/1,712 recalled instances. Every ceiling uses exact integer optimization. The 1,000 archived random-ranking draws have $K=20$ 2.5/97.5 percentiles of 0.077/0.104.

\subsection*{Validation-queue selection and allocation-policy analysis}

The queue analysis selected logistic regularization by validation $K=20$ one-to-one qualified recall, followed by many-to-one recall, precision, and $C$ closest to 1 on the log scale. The selected value was $C=10$. The model was refitted on train plus validation before test scoring. The post-test 5--50 policy maximized total score under a 520-candidate budget and is reported as exploratory. Unconstrained global top-520 is a diagnostic upper bound. The validation threshold is the boundary score of the validation global top-520 queue; transferring that fixed score to test accepted 426 candidates.

Candidate-level score discrimination was heterogeneous before allocation. Pooled AUC/AP were 0.9605/0.8617, whereas date-macro values were 0.8774/0.6901. Table~\ref{tab:dateauc} reports each date. Twenty-four of 26 tiles contained both classes and supported tile-level AUC/AP; their macro values were 0.9091/0.6869, with worst-tile values of 0.6667/0.0263. These values are descriptive and are not treated as independent population replicates.

\begin{table}[htbp]
\centering
\caption{Candidate-level role-qualified discrimination for the queue-selected ranker by acquisition date. Class counts accompany each AUC/AP value.}
\label{tab:dateauc}
\small
\begin{tabular}{lrrrrr}
\toprule
Stratum & Candidates & Weed positive & Other & AUC & Average precision \\
\midrule
25 May & 822 & 15 & 807 & 0.8798 & 0.2131 \\
30 May & 955 & 132 & 823 & 0.9535 & 0.7492 \\
6 June & 517 & 279 & 238 & 0.9381 & 0.9402 \\
15 June & 57 & 38 & 19 & 0.7382 & 0.8577 \\
\midrule
Pooled & 2,351 & 464 & 1,887 & 0.9605 & 0.8617 \\
Date macro & -- & -- & -- & 0.8774 & 0.6901 \\
\bottomrule
\end{tabular}
\end{table}

\begin{table}[htbp]
\centering
\caption{Queue allocation using the validation-queue-selected masked-DINOv2 ranker. The post-test 5--50/tile and unconstrained global top-$N$ policies are exploratory diagnostics.}
\scriptsize
\resizebox{\textwidth}{!}{%
\begin{tabular}{lrrrrrr}
\toprule
Policy & Accepted & Precision & Spatial & Qualified & 1:1 qualified & Any crop \\
\midrule
Fixed $K=20$ & 520 & 275/520 & 390/1712 & 324/1712 & 227/1712 & 13/26 \\
Exploratory 5--50/tile & 520 & 382/520 & 533/1712 & 453/1712 & 315/1712 & 14/26 \\
Diagnostic global top-$N$ & 520 & 388/520 & 547/1712 & 464/1712 & 320/1712 & 12/26 \\
Validation threshold & 426 & 348/426 & 485/1712 & 424/1712 & 290/1712 & 8/26 \\
\bottomrule
\end{tabular}}
\end{table}

\begin{table}[htbp]
\centering
\caption{Per-date one-to-one qualified recall and accepted candidates for fixed, exploratory 5--50, and diagnostic global allocation.}
\scriptsize
\resizebox{\textwidth}{!}{%
\begin{tabular}{lrrrrrr}
\toprule
Date & Fixed $N$ & Fixed recall & Bounded $N$ & Bounded recall & Global $N$ & Global recall \\
\midrule
25 May & 160 & 10/392 (0.0255) & 51 & 6/392 (0.0153) & 38 & 6/392 (0.0153) \\
30 May & 160 & 71/440 (0.1614) & 149 & 77/440 (0.1750) & 153 & 78/440 (0.1773) \\
6 June & 160 & 135/726 (0.1860) & 285 & 221/726 (0.3044) & 293 & 225/726 (0.3099) \\
15 June & 40 & 11/154 (0.0714) & 35 & 11/154 (0.0714) & 36 & 11/154 (0.0714) \\
\bottomrule
\end{tabular}}
\end{table}

For global minus fixed one-to-one qualified recall, the 26 paired tile differences had median 0, interquartile range 0--0.0763, and range $-0.1176$ to 0.2027. Nine tiles improved, 12 were unchanged, and 5 declined. Many-to-one differences had median 0, interquartile range 0--0.1000, and range $-0.1176$ to 0.3919, with the same 9/12/5 counts. Two capture-index proxy groups yielded global qualified recall of 0.2574 and 0.2806; corresponding fixed values were 0.1782 and 0.1970.

\begin{table}[htbp]
\centering
\caption{Allocation concentration across 26 test tiles. The 6 June share is the fraction of each policy's accepted candidates assigned to that acquisition date.}
\small
\setlength{\tabcolsep}{4pt}
\begin{tabular}{lrrrrr}
\toprule
Policy & Min/tile & Zero tiles & Tiles $<5$ & Tiles $<10$ & Max/tile \\
\midrule
Fixed $K=20$ & 20 & 0 & 0 & 0 & 20 \\
Exploratory 5--50 & 5 & 0 & 0 & 8 & 50 \\
Diagnostic global & 0 & 1 & 5 & 8 & 55 \\
Validation threshold & 0 & 2 & 7 & 11 & 45 \\
\bottomrule
\end{tabular}

\vspace{0.5em}
\begin{tabular}{lrr}
\toprule
Policy & Allocation Gini & 6 June share \\
\midrule
Fixed $K=20$ & 0.0000 & 0.3077 \\
Exploratory 5--50 & 0.3817 & 0.5481 \\
Diagnostic global & 0.4220 & 0.5635 \\
Validation threshold & 0.4648 & 0.5986 \\
\bottomrule
\end{tabular}
\end{table}

\begin{table}[htbp]
\centering
\caption{Date-macro and worst-date queue metrics. Worst denotes the minimum across the four dates.}
\small
\setlength{\tabcolsep}{4pt}
\begin{tabular}{p{3.8cm}rrrrr}
\toprule
Policy & Macro P & Macro Q & Macro 1:1 & Worst Q & Worst 1:1 \\
\midrule
Fixed $K=20$ & 0.5516 & 0.1504 & 0.1111 & 0.0255 & 0.0255 \\
Exploratory 5--50 & 0.6133 & 0.1959 & 0.1415 & 0.0153 & 0.0153 \\
Diagnostic global & 0.6160 & 0.2001 & 0.1435 & 0.0153 & 0.0153 \\
Validation threshold & 0.6876 & 0.1836 & 0.1311 & 0.0128 & 0.0128 \\
\bottomrule
\end{tabular}
\end{table}

\begin{table}[htbp]
\centering
\caption{Crop exposure severity. Crop-instance rows count 256 official-test crop instances.}
\scriptsize
\resizebox{\textwidth}{!}{%
\begin{tabular}{lrrrrrr}
\toprule
Policy & Crop-overlap cand. & Crop frac. 95th & Crop frac. max & Crop inst. $\ge0.10$ & $\ge0.25$ & $\ge0.50$ \\
\midrule
Fixed $K=20$ & 57/520 & 0.5910 & 0.9730 & 52/256 & 50/256 & 32/256 \\
Diagnostic global top-$N$ & 26/520 & 0.0036 & 0.9583 & 21/256 & 19/256 & 10/256 \\
Exploratory 5--50 & 28/520 & 0.1148 & 0.9189 & 23/256 & 20/256 & 12/256 \\
Validation threshold & 14/426 & 0 & 0.9189 & 12/256 & 10/256 & 5/256 \\
\bottomrule
\end{tabular}}
\end{table}

\subsection*{Equal-target-mask-access semantic baseline}

The archived 40-epoch distributed reference run followed its specified split, optimizer, augmentation, threshold grid, area grid, and selection rule. Validation selected epoch 36, threshold 0.30, and 16-pixel minimum area. Validation three-class/plant-role mean IoU were 0.7637/0.6582; fixed-$K$ one-to-one qualified recall was 0.1908. Official-test semantic IoU was 0.9778 background, 0.6957 crop, and 0.7131 weed (three-class mean 0.7955; plant-role mean 0.7044). A separate five-seed H200 series quantifies training variability under the same single-device recipe.

\begin{table}[htbp]
\centering
\caption{Official-test target-semantic queue. Components are ranked by their target-semantic weed score; validation alone selected epoch and component settings.}
\small
\resizebox{\textwidth}{!}{%
\begin{tabular}{lrrrrrrrr}
\toprule
Model & Candidates & Accepted & Weed cand. & Precision & Spatial & Qualified & 1:1 qualified & Crop scenes \\
\midrule
Target semantic, fixed $K=20$ & 1,347 & 518 & 355 & 0.6853 & 520/1712 (0.3037) & 453/1712 (0.2646) & 348/1712 (0.2033) & 6/26 (0.2308) \\
\bottomrule
\end{tabular}}
\end{table}

\begin{table}[htbp]
\centering
\caption{Complete representation--allocation comparison. Exploratory 5--50 policies preserve the fixed total within each representation.}
\footnotesize
\setlength{\tabcolsep}{2pt}
\begin{tabular}{L{3.0cm}L{2.3cm}rrrrr}
\toprule
Representation & Allocation & $N$ & Prec. & 1:1 & Crop & Worst 1:1 \\
\midrule
Masked DINO & Fixed 20/tile & 520 & 0.5288 & 0.1326 & 13/26 & 0.0255 \\
Masked DINO & Exploratory 5--50 & 520 & 0.7346 & 0.1840 & 14/26 & 0.0153 \\
Target semantic & Fixed 20/tile & 518 & 0.6853 & 0.2033 & 6/26 & 0.1169 \\
Target semantic & Exploratory 5--50 & 518 & 0.8205 & 0.2477 & 1/26 & 0.0510 \\
\bottomrule
\end{tabular}
\end{table}

\subsection*{Five-seed H200 stability series}

Five independently initialized single-device runs used the same 40-epoch recipe, official split, validation-only selection rule, and batch size of six. All runs completed on NVIDIA H200 devices. Validation selected the 16-pixel component minimum in every run; the foreground threshold varied across 0.30, 0.50, and 0.70. Table~\ref{tab:supp_h200_seeds} reports every seed rather than treating candidates, tiles, or plant instances as independent training repetitions.

\begin{table}[htbp]
\centering
\caption{Per-seed target-semantic H200 results. Seed values are numeric identifiers, not acquisition dates. Fixed precision and one-to-one recall use $K=20$ per tile. The exploratory 5--50 allocation preserves the fixed accepted total within each seed.}
\label{tab:supp_h200_seeds}
\footnotesize
\setlength{\tabcolsep}{3pt}
\resizebox{\textwidth}{!}{%
\begin{tabular}{rrrrrrrrr}
\toprule
Seed ID & Epoch & $T$ & $A$ & Semantic mIoU & Fixed prec. & Fixed 1:1 & 5--50 prec. & 5--50 1:1 \\
\midrule
20260810 & 38 & 0.5 & 16 & 0.7931 & 0.7112 & 0.2062 & 0.8841 & 0.2593 \\
20260811 & 36 & 0.7 & 16 & 0.7790 & 0.7248 & 0.1992 & 0.9188 & 0.2605 \\
20260812 & 40 & 0.3 & 16 & 0.8116 & 0.6879 & 0.2033 & 0.8786 & 0.2605 \\
20260813 & 39 & 0.5 & 16 & 0.7897 & 0.7569 & 0.2097 & 0.9216 & 0.2664 \\
20260814 & 31 & 0.5 & 16 & 0.8041 & 0.7035 & 0.1951 & 0.8663 & 0.2424 \\
\bottomrule
\end{tabular}}
\end{table}

The fixed queue accepted $511.8\pm5.6$ candidates across seeds (mean $\pm$ sample SD). Fixed precision and one-to-one recall were $0.7168\pm0.0260$ and $0.2027\pm0.0057$; exploratory 5--50 allocation increased them to $0.8939\pm0.0249$ and $0.2578\pm0.0090$. The paired 5--50-minus-fixed differences were $+0.1770\pm0.0146$ for precision and $+0.0551\pm0.0053$ for one-to-one recall, and both were positive in all five seeds. Crop-exposed scene frequency decreased by $0.1846\pm0.0632$ in all five seeds. Worst-date one-to-one recall decreased by $0.0627\pm0.0070$ in all five seeds.

For the stability analysis, each selected checkpoint was replayed on the official test images to regenerate component masks and candidate--instance edges. Candidate IDs and fixed-queue precision, many-to-one recall, one-to-one recall, and crop exposure were checked against the archived run output for every seed and every test tile before alternative allocation policies were evaluated. The reference replay regenerated all 1,347 candidate identifiers and reproduced all 26 fixed-queue scene records. The archive also preserves an unsuccessful box-only reconstruction that overcounted one recovered instance in one seed.

\begin{table}[htbp]
\centering
\caption{Complete post-test quota grid over the five H200 seeds. Each cell reports mean pooled / mean worst-date one-to-one role-qualified recall. ``None'' denotes no maximum.}
\label{tab:supp_quota_grid}
\small
\setlength{\tabcolsep}{5pt}
\resizebox{\textwidth}{!}{%
\begin{tabular}{lrrrr}
\toprule
Minimum & Maximum 20 & Maximum 30 & Maximum 50 & Maximum none \\
\midrule
0  & 0.2027 / 0.1100 & 0.2484 / 0.0558 & 0.2643 / 0.0243 & 0.2643 / 0.0210 \\
5  & 0.2027 / 0.1100 & 0.2468 / 0.0583 & 0.2578 / 0.0473 & 0.2578 / 0.0473 \\
10 & 0.2027 / 0.1100 & 0.2405 / 0.0809 & 0.2442 / 0.0783 & 0.2439 / 0.0783 \\
20 & 0.2027 / 0.1100 & 0.2027 / 0.1100 & 0.2027 / 0.1100 & 0.2027 / 0.1100 \\
\bottomrule
\end{tabular}}
\end{table}

Validation policy selection required validation candidate-level role-qualified precision at least as high as fixed $K=20$, then maximized validation worst-date one-to-one recall, pooled one-to-one recall, and precision. It selected 0--20 in seed IDs 20260810, 20260812, 20260813, and 20260814, and 0--30 in seed ID 20260811. Thus, the exploratory 5--50 setting was not selected by the group-aware validation rule. Date-balanced global allocation yielded precision $0.7710\pm0.0226$, pooled one-to-one recall $0.2192\pm0.0064$, and worst-date recall $0.1163\pm0.0190$. Its paired one-to-one effect relative to fixed allocation was negative in all five seeds for capture-index patch proxy 0 (mean $-0.0458$) and positive in all five for proxy 1 (mean $+0.0603$). The 5--50 effect was positive in both proxies for all five seeds.

IoU-only weed-instance recall at thresholds 0.25/0.50/0.75 was 429/1,712 (0.2506), 280/1,712 (0.1636), and 53/1,712 (0.0310). The median, 75th percentile, and 95th percentile queued crop fraction were all zero; the maximum was 0.9474. Two tiles contained fewer than 20 components, so the accepted total was 518 rather than 520.

\begin{table}[htbp]
\centering
\caption{Target semantic fixed-$K$ queue by acquisition date. Every proportion is accompanied by an archived per-tile integer count.}
\small
\resizebox{\textwidth}{!}{%
\begin{tabular}{lrrrrrrrr}
\toprule
Date & Tiles & Candidates & Accepted & Precision & Spatial & Qualified & 1:1 qualified & Crop scenes \\
\midrule
25 May & 8 & 304 & 158 & 70/158 (0.4430) & 88/392 (0.2245) & 68/392 (0.1735) & 63/392 (0.1607) & 6/8 \\
30 May & 8 & 446 & 160 & 124/160 (0.7750) & 175/440 (0.3977) & 144/440 (0.3273) & 124/440 (0.2818) & 0/8 \\
6 June & 8 & 534 & 160 & 143/160 (0.8938) & 226/726 (0.3113) & 211/726 (0.2906) & 143/726 (0.1970) & 0/8 \\
15 June & 2 & 63 & 40 & 18/40 (0.4500) & 31/154 (0.2013) & 30/154 (0.1948) & 18/154 (0.1169) & 0/2 \\
\bottomrule
\end{tabular}}
\end{table}

Across the four dates, target-semantic macro precision, many-to-one spatial recall, many-to-one qualified recall, and one-to-one qualified recall were 0.6404, 0.2837, 0.2465, and 0.1891. The corresponding worst-date values were 0.4430, 0.2013, 0.1735, and 0.1169. Crop-overlap scene frequency had a date macro of 0.1875 and a worst-date value of 0.7500. The repeated H200 series evaluates the same metric family across independent initializations.

The date pattern differs from the commissioned masked-DINO queue. Target semantic training improves the 25 and 30 May fixed queues, while the two-image 15 June stratum reaches 0.1169 one-to-one recall. These date strata describe temporal variation within the study field.

\section{S6. Matching-contract, IoU, crop-exposure, and multiplicity sensitivity}

The evaluated queue-contract table contains every combination of model, $K\in\{1,5,10,20,50\}$, date/pooled stratum, instance coverage $\{0.25,0.50,0.75\}$, labeled coverage $\{0.25,0.50,0.75\}$, and purity $\{0.50,0.75,0.90\}$. For component-masked DINOv2 at $K=20$, pooled many-to-one qualified recall ranges 0.0485--0.2886, and one-to-one recall ranges 0.0374--0.1881. At the main 0.50/0.50/0.90 contract, these values are 319/1,712 (0.1863) and 230/1,712 (0.1343). Under IoU-only matching at thresholds 0.25/0.50/0.75, many-to-one recall is 325/1,712, 163/1,712, and 24/1,712; one-to-one recall is 307/1,712, 163/1,712, and 24/1,712.

\begin{table}[htbp]
\centering
\caption{Crop-overlap scene frequency at $K=20$. A scene is counted when at least one queued region has crop fraction greater than or equal to the listed threshold; for threshold zero, positive overlap is required in implementation.}
\small
\resizebox{\textwidth}{!}{%
\begin{tabular}{lrrrr}
\toprule
Model & Any $>0$ & $\ge0.01$ & $\ge0.10$ & $\ge0.50$ \\
\midrule
Frozen DINOv2 & 10/26 & 10/26 & 10/26 & 9/26 \\
Masked DINOv2 & 11/26 & 11/26 & 11/26 & 9/26 \\
\bottomrule
\end{tabular}}
\end{table}

The archive additionally reports queued candidate--instance multiplicity and IoU matching at 0.25, 0.50, and 0.75. These tables quantify candidates covering several small weeds and weeds fragmented across several candidates, complementing the main many-to-one metric.

\section{S7. Restricted training-label draws with fully labelled validation and fine-tuning}

\begin{table}[htbp]
\centering
\caption{Mean official-test results over ten draws of 300 simulated training-candidate labels derived from dense masks. All models also use the fully labelled validation candidate set.}
\small
\resizebox{\textwidth}{!}{%
\begin{tabular}{lrrrr}
\toprule
Selection & AUC & AP & $K=20$ precision & $K=20$ qualified recall \\
\midrule
Prospective uniform & 0.9549 & 0.8544 & 0.5242 & 0.1724 \\
Prospective DINOv2 coreset & 0.9543 & 0.8534 & 0.5217 & 0.1717 \\
Retrospective stratified & 0.9547 & 0.8502 & 0.5231 & 0.1725 \\
\bottomrule
\end{tabular}}
\end{table}

All ten candidate-ID lists and their SHA-256 checksums are supplied in the reproducibility archive. Coreset selection is greedy farthest-first in frozen DINOv2 embedding space with a seeded initial point; uniform selection is seeded sampling without replacement. Training data determine candidate IDs, and official validation selects $C$ and epoch. The 2.5/97.5 draw percentiles reported below are descriptive empirical quantiles, not confidence intervals.

\begin{table}[htbp]
\centering
\caption{Paired last-block fine-tuning differences (fine-tuned minus frozen) on identical coreset draws.}
\small
\begin{tabular}{lrrr}
\toprule
Metric & Mean difference & Empirical 2.5th & Empirical 97.5th \\
\midrule
AUC & $-0.0396$ & $-0.0525$ & $-0.0170$ \\
$K=20$ precision & $-0.0373$ & $-0.0476$ & $-0.0169$ \\
$K=20$ qualified recall & $-0.0115$ & $-0.0234$ & $-0.0001$ \\
\bottomrule
\end{tabular}
\end{table}

\begin{table}[htbp]
\centering
\caption{Per-repeat paired AUC differences for last-block fine-tuning minus the frozen representation. Repeats use identical selected training IDs within each pair.}
\small
\resizebox{\textwidth}{!}{%
\begin{tabular}{rrrrrrrrrrr}
\toprule
Repeat & 0 & 1 & 2 & 3 & 4 & 5 & 6 & 7 & 8 & 9 \\
\midrule
$\Delta$AUC & $-0.0508$ & $-0.0367$ & $-0.0412$ & $-0.0530$ & $-0.0424$ & $-0.0456$ & $-0.0116$ & $-0.0353$ & $-0.0400$ & $-0.0392$ \\
\bottomrule
\end{tabular}}
\end{table}

The predefined fine-tuning configuration used DINOv2 ViT-S/14, tight unpadded crops, 8 epochs, batch 64, AdamW, and class-weighted binary cross-entropy. Weight decay was $10^{-4}$, the last-block/norm learning rate was $10^{-5}$, and the head learning rate was $10^{-3}$. Horizontal and vertical flips supplied augmentation. The final block, final normalization, and head were trainable, while earlier blocks remained frozen. Validation selected epoch 8 in every repeat. The paired estimates characterize this archived recipe.

\section{S8. Statistical interpretation}

The main tables report raw test numerators and denominators, pooled values, all four acquisition dates, and paired tile differences. The 15 June date contains two official-test images and consequently receives greater relative weight in an unweighted date macro. The archive supplies both date-macro and tile/instance-weighted summaries. Global-versus-fixed comparisons include the paired median, interquartile range, full range, and improved/unchanged/worsened counts. These statistics describe the official spatial split.

The release supplies capture indices as its spatial coordinate proxy. A full-tile DINOv2 Hungarian linkage to nominal locations produced negative assignment margins for 12 of 18 test assignments. We therefore used the largest ordered-index gap to form two 13-image capture-index proxy groups. Their bootstrap summaries quantify sensitivity to this grouping. Repeat-wise 2.5th and 97.5th percentiles are empirical simulation quantiles.

\section{S9. Development diagnostics}

\begin{table}[htbp]
\centering
\caption{Date-grouped development diagnostics. The official spatial analysis is reported separately.}
\small
\resizebox{\textwidth}{!}{%
\begin{tabular}{p{3.1cm}p{4.5cm}p{6.1cm}}
\toprule
Analysis & Result & Interpretation \\
\midrule
Nested Platt calibration & Date-macro Brier 0.2112 before, 0.2776 after; NLL 0.6338 before, 0.8239 after & Date transfer increased both calibration errors; queue analyses use ranking scores \\
Empirical threshold controller & Accepted precision 0.1317, empty response 0.7448, crop-scene frequency 0.1042 & Training-date threshold transfer produced a sparse held-out queue with low precision \\
Cost-sensitive ranker, $K=20$ & AUC 0.7301, AP 0.7272, precision 0.4327, instance recall 0.0842, crop-scene frequency 0.3542 & Crop weighting shifted recovery and exposure in opposite directions \\
Role $\times$ support, $K=20$ & AUC 0.7151 versus 0.8542 role-only; recall 0.0910 versus 0.0520; crop frequency 0.3698 versus 0.1510 & Support multiplication raised recovery and crop exposure while weakening ranking \\
Last-block DINOv2 fine-tuning & Mean paired AUC $-0.0396$ and qualified recall $-0.0115$ & Small target supervision favored frozen representations \\
\bottomrule
\end{tabular}}
\end{table}

The run registry retains every diagnostic, including unsuccessful executions, under its original identifier. Crop-overlap frequency is defined consistently across these records.

\section{S10. Runtime, environment, and model provenance}

\begin{table}[htbp]
\centering
\caption{Principal runtime and memory records.}
\small
\resizebox{\textwidth}{!}{%
\begin{tabular}{lrrl}
\toprule
Component & Wall time & Peak CUDA memory & Device \\
\midrule
Unmasked DINOv2 embeddings & 93.94 s & 601.9 MB & H200, physical GPU 0 \\
Masked DINOv2 embeddings & 32.52 s & 601.9 MB & H200, physical GPU 0 \\
Official proposal audit & 35.46 s & CPU & -- \\
Queue-policy and tile-level audit & 7.63 s & CPU & -- \\
Target semantic training, validation, and finalization & 2.86 h & See run history & H200, physical GPUs 5--7 \\
Target semantic test inference & 0.931 s (0.0358 s/tile) & 583.8 MB & H200, physical GPU 5 \\
Five-seed target-semantic training & $712.48\pm8.22$ s/run & See run history & H200, one device/run on physical GPUs 6--7 \\
Five-seed target-semantic test replay & $0.950\pm0.021$ s/run & 513.4 MiB/run & H200, one device/run on physical GPUs 6--7 \\
Ten fine-tuning repeats & 93.89 s elapsed & 660.1 MB/repeat & H200, physical GPUs 0--7 \\
\bottomrule
\end{tabular}}
\end{table}

\begin{table}[htbp]
\centering
\caption{Supervision and computation for the two representation paths. Runtime entries are measured records rather than hardware-normalized estimates.}
\footnotesize
\resizebox{\textwidth}{!}{%
\begin{tabular}{L{3.1cm}L{3.0cm}L{3.4cm}L{3.5cm}L{3.0cm}}
\toprule
Path & Target-mask access & Additional labelled data & Public pretraining & Recorded computation \\
\midrule
Source specialists $+$ masked DINO & 104 train masks for proposal selection and 9,887 candidate labels; 26 validation masks and 2,436 candidate labels for selection & 7,050 SugarBeets training masks; 901 source-validation masks & ImageNet-1K ResNet-18; DINOv2 ViT-S/14 LVD-142M & Masked test embeddings: 32.52 s, 601.9 MB; source training time was not retained \\
Target semantic & 104 train masks; 26 validation masks for epoch and component selection & None & ImageNet-1K ResNet-18 & Training, validation, and finalization: 2.86 h on three H200s; test: 0.931 s, 583.8 MB \\
Target-semantic stability series & Same 104 train and 26 validation masks in every seed & None & ImageNet-1K ResNet-18 & Five single-device 40-epoch runs: $712.48\pm8.22$ s/run; exact test replay: $0.950\pm0.021$ s/run and 513.4 MiB \\
\bottomrule
\end{tabular}}
\end{table}

The 2.86-h distributed reference and the approximately 712-s single-device repetitions are records from different implementations and are not a hardware speed benchmark. The reference used a three-process distributed path with local batch size two and performed full validation setting search and finalization during its 40-epoch run. The later repetitions used a single-process path with local batch size six and were launched concurrently across two H200 accelerators. Model architecture, global batch size, optimization objective, epoch count, and validation selection contract were held constant; process topology, validation orchestration, and logging overhead were not. The archive therefore reports both measurements without converting them into a speedup claim.

The recorded GPU environment is Python 3.10.12, PyTorch 2.7.1+cu126, timm 1.0.28, and NumPy 2.2.6. The public DINOv2 ViT-S/14 LVD-142M weights were obtained through timm. Model, embedding, configuration, and output checksums are recorded in the archive manifest.

\section{S11. Reproducibility and archive record}

The versioned archive is available at \url{https://github.com/Fytap/agrispec-vlm-reproducibility/releases/tag/v1.1.0}. It preserves the executable directory structure, frozen configurations, declared inputs, source-data tables, environment specification, run manifests, and unsuccessful execution records. Clean-room commands recompute tables and figures from archived predictions and exact matching outputs. Upstream source files, public-weight revisions, checkpoint hashes, acquisition instructions, and recorded training metadata document model construction. The release adds the complete proposal grid, five H200 checkpoints, exact component-mask replay, quota and spatial-group analyses, and verification hashes.

\begingroup
\footnotesize
\setlength{\tabcolsep}{3pt}
\renewcommand{\arraystretch}{0.94}
\begin{longtable}{>{\raggedright\arraybackslash}p{2.5cm}>{\raggedright\arraybackslash}p{3.8cm}>{\raggedright\arraybackslash}p{6.5cm}}
\caption{Computational provenance and audit controls.}\\
\toprule
Element & Verification & Archived record \\
\midrule
\endfirsthead
\toprule
Element & Verification & Archived record \\
\midrule
\endhead
Metric layers & Separate spatial, role-qualified, and queued recall & Definitions, thresholds, and integer counts \\
Split/merge & Sensitivity at 0.10, 0.25, and 0.50 instance overlap & Threshold grid and per-setting summaries \\
Scene denominator & Crop-exposure frequencies use all 26 test tiles & Tile-level queue and overlap records \\
Spatial grouping & Released capture indices define the spatial proxy & Date, patch, tile, and capture-index summaries \\
Candidate support & Rankers use the complete commissioned pool & Candidate predictions and support indicators \\
DINOv2 context & Paired unmasked and component-masked crops & Embedding manifests and paired queue outputs \\
Fine-tuning & Ten paired frozen versus last-block repeats & Seeds, checkpoints, logs, and paired differences \\
Target semantic model & Archived split, selection rule, and hyperparameters & Successful and unsuccessful execution records \\
Target semantic allocation & Fixed, complete quota-grid, date-balanced, and validation-selected policies & Per-scene selections, exact component-mask intersections, and hashes \\
H200 seed stability & Five independent initializations with identical data roles and optimization recipe & Checkpoints, hardware attestations, per-seed summaries, exact-mask policy replay, paired differences, and SHA-256 manifest \\
Raster accounting & Released-mask and paper-reported totals & Consistency and exclusion audits \\
\bottomrule
\end{longtable}
\endgroup

\end{document}
